\documentclass[11pt,a4paper]{article}
\usepackage[T1]{fontenc}
\usepackage{amssymb}
\usepackage{isabelle,isabellesym}
\usepackage{railsetup}
\usepackage{mathpartir}
\usepackage[skip=0.5ex]{parskip}
\usepackage{enumitem}
\usepackage{xcolor}
% a described form on a line of its own, its explanation indented below
% and a thin rule after a list of such forms, before the next diagram
\newif\ifdescformseen
\setlist[description]{leftmargin=1.5em, labelindent=0pt, labelwidth=0pt, labelsep=0pt,
  itemsep=0.6ex, after=\descformrule}
\newcommand{\descformrule}{\ifdescformseen\global\descformseenfalse
  \endgraf\vspace{1ex}\noindent{\color{lightgray}\rule{\linewidth}{0.4pt}}\endgraf\fi}
\newcommand{\descform}[1]{\global\descformseentrue\hspace*{-\leftmargin}#1\\*}
\setlist[itemize]{leftmargin=1.2em}

% room between a syntax diagram and the text below it; underscores in rule names
\AddToHook{env/railoutput/after}{\vspace{1.5ex}}
\makeatletter
\def\rail@nontfont{\isabellestyle{literalunderscore}}
\def\rail@namefont{\isabellestyle{literalunderscore}}
\makeatother

\usepackage{pdfsetup}

\urlstyle{rm}
\isabellestyle{literalunderscore}

\isadroptag{theory}

% show the quotes of strings: the manual quotes only literal spellings
\renewcommand{\isachardoublequoteopen}{\textquotedbl}
\renewcommand{\isachardoublequoteclose}{\textquotedbl}
\renewcommand{\isachardoublequote}{\textquotedbl}


\begin{document}

\title{Tokenization, Indentation-sensitive Grammars and Executable Recognizers}
\author{Basil F\"urer \and Dmitriy Traytel}
\maketitle

\begin{abstract}
  Two definitional packages for Isabelle/HOL that turn the source text of a
  language into verified values: The command \isakeyword{lexer} elaborates a
  flex-style specification, with modes, longest match and rule priority, into an
  inductive lexing judgment and an executable lexer on precomputed DFA tables.
  The command \isakeyword{grammar} elaborates an indentation-sensitive context-free
  grammar with semantic actions into a grammar specification, decoders from parse
  trees to the values of the semantic actions and an inductive relation. The related
  command \isakeyword{parser} emits an executable LR parser for it.
  Both executable constants for lexing and parsing are proved sound and complete
  with respect to their specifications. The parser completeness, and with it the
  unambiguity of the grammar, follows from an LALR conflict analysis that the
  command carries out by evaluation.
\end{abstract}

\newpage

\tableofcontents

\newpage

\input{TIGER.tex}

\newpage

\bibliographystyle{abbrv}
\bibliography{root}

\end{document}
